\chapter{Conclusions}
This project set out to investigate whether the shape of a dielectric resonator
can be optimised to enhance the Purcell factor, and consequently improve the
performance of a solid-state MASER operating at room temperature and zero
magnetic field. Using COMSOL Multiphysics\textregistered{}, COMSOL API and Java,
a genetic algorithm (GA) was implemented to iteratively optimise the geometry of
a \ce{SrTiO3}-based resonator. The primary objective was to maximise the ratio
of quality factor ($Q_f$) to mode volume ($V_m$), which are the main metrics
that determine the Purcell factor of the resonator and is very important in
determining the efficiency of microwave amplification.

The first significant outcome of this work was the successful creation and
validation of an initialised geometry that exhibited strong
$\text{TE}_{01\delta}$ mode confinement at the target frequency of 1.45 GHz.
This was achieved by modifying the original MASER-in-a-Shoebox dimensions and
tuning the resonator using geometric adjustments, mimicking the real-world
tuning process via adjustable cavity ceilings. This model provided a baseline
for optimisation and ensured that all GA-evaluated geometries could be fairly
compared.

The GA produced a series of candidate geometries, with the best-performing
structure demonstrating a marked improvement in Purcell factor compared to the
baseline model. While this improvement was obtained under idealised (lossless)
conditions, it confirmed that purely geometrical changes can enhance the
electromagnetic confinement properties of the resonator. Julia-based simulations
using \texttt{QuantumCumulants.jl} showed that this translated into increased
output power under strong coupling assumptions, although these results are
illustrative rather than predictive, serving as a visual tool to compare
performance between the baseline and optimised models. However, when the
dielectric loss of \ce{SrTiO3} was introduced into the simulation, the GA was
very slow at discovering more optimal configurations and exhibited mode
switching. Thus making any results invalid. After existed that configurations
with a higher Purcell factor existed, the results were further examined to
determine where the root of the error was. It was ultimately concluded that the
code contained a bug.

This project has shown that resonator geometry can play a meaningful role in
enhancing MASER performance. Despite the results be preliminary, they (and the
code produced) should hopefully serve as a framework for future investigations
into this area of research. The Java-based GA proved a suitable prototype
however more sophisticated numerical control is required as Java lacks the
native libraries and numerical tools. Hence, the migration to COMSOL’s
LiveLink\texttrademark{} for MATLAB\textregistered{} environment has been
proposed. This environment can utilise the more flexible and numerically robust
to implement new optimisation techniques, such as machine learning or
simulations in 3D.

In summary, this project has demonstrated that resonator geometry significantly
affects key MASER performance metrics. While lossless simulations provide a
useful sandbox for testing ideas, realistic modelling and more advanced
optimisation tools will be essential for translating these insights into
deployable MASER technologies.
